\chapter*{Abstract}
\addcontentsline{toc}{chapter}{Abstract}

Security Operations Centres (SOCs) currently face two mutually reinforcing problems. First,
the volume of alerts produced by detection tooling far exceeds human analysis capacity,
leading to the well-documented phenomenon of \textit{alert fatigue}: analysts dismiss or
postpone potentially relevant warnings simply because there is no time to review them all.
Second, those alerts arrive \textit{without context}: an isolated warning does not indicate
whether the observed behaviour corresponds to a documented technique, whether it has been
used by a known threat actor, or which organisational assets are consequently at risk.

Compounding this, targeted attacks do not manifest as isolated events but as
\textbf{multi-stage sequences} unfolding over time: reconnaissance, initial access,
execution, persistence, lateral movement and exfiltration. Detection systems that evaluate
each event independently are, by construction, unable to recognise that temporal structure,
and are therefore easily evaded through individually unremarkable techniques.

This Bachelor's Thesis proposes the design, implementation and validation of a data platform
addressing both problems jointly. The system ingests heterogeneous network and host
telemetry, normalises it to a canonical schema (OCSF), detects multi-stage attack sequences
via complex event processing (CEP), contextualises them against a cyber threat intelligence
knowledge graph aligned with MITRE ATT\&CK, reasons over feasible attack paths towards
critical assets using a Datalog-based logic engine, and finally produces an
\textbf{explainable} incident prioritisation in which every score is accompanied by a
breakdown of the contributing factors.

The architecture rests on three pillars: multi-stage detection, the reasoned attack graph,
and explainable prioritisation. The ultimate goal is not to replace existing detection
tooling, but to transform its output into actionable information, providing
\textbf{cyber situational awareness} and decision support to the analyst.

Validation is carried out on public datasets labelled with MITRE ATT\&CK techniques,
enabling a quantitative assessment of detection accuracy, alert volume reduction and
prioritisation quality against established baselines.

\vspace{1em}
\noindent\textbf{Keywords:} cyber situational awareness, intrusion detection, MITRE ATT\&CK,
complex event processing, attack graphs, threat intelligence, real-time data engineering,
explainability.
